\section{Non-LLM Psychometrics}

\subsection{Goal}

The goal of the non-LLM workstream was not just to collect alternative ARC scores. It was to ask a sharper question: once we add bona fide non-LLM systems to the repo, do they look more human-like than the LLM aggregate, less human-like, or simply different in a complementary way?

\subsection{Data}

The main comparison space is ARC-AGI-2 Public Eval because that is where the human attempt data and stored prediction corpora overlap cleanly. The primary subset contains 110 public-eval test pairs with at least 8 human attempts each. The main non-LLM sources are TRM ARC-AGI-II evaluator submissions and VARC ARC-2 prediction dumps. CompressARC contributes an ARC-1 sidecar rather than a direct ARC-2 comparison.

\subsection{Methods}

The non-LLM section uses more than one null because low-score binary systems are easy to misread. It benchmarks systems against the human split-half distribution, reruns the analysis at several human-attempt thresholds, tests fixed-accuracy random-placement nulls, controls for simple task features, compares against accuracy-matched weak LLMs, and checks for residual signal after controlling the LLM average. It also looks for complementarity by asking whether non-LLM systems rescue human-easy and LLM-hard items above chance.

\subsection{Experiment}

The core experiment compares several selected profiles on the same ARC-2 human-overlap subset: the LLM average, the best-aligned single LLM, the best-score single LLM, a mid-training TRM checkpoint, a later higher-score TRM checkpoint, and the strongest stored VARC ARC-2 profile. The follow-up experiments then ask whether low raw correlations are better than score-matched random placement, whether any residual human-like structure survives simple controls, and whether TRM training improves human-likeness or only score.

\subsection{Results}

The top line is disappointing but clear. The LLM average reaches a human-correlation value around 0.402 on the primary 110-pair subset, while the best-aligned single LLM reaches about 0.439. The most human-aligned TRM checkpoint reaches only about 0.214, and the best VARC profile about 0.158. Human-equivalence is not rejected for the LLM average or the best-aligned LLM, but it is rejected for the leading non-LLM profiles.

At the same time, raw correlation is not the whole story. The threshold-sensitivity table shows that the ordering is stable as the human-attempt filter tightens: the LLM aggregate and the best-aligned single LLM improve steadily, while TRM and VARC remain well below them. The mid-training TRM checkpoint exceeds the same-accuracy random-placement null, whereas the later best-score TRM checkpoint does not. TRM and VARC also rescue a handful of human-easy and LLM-hard items, and the TRM plus VARC union rescues 6 such items above chance. So the non-LLM systems are not useless or redundant; they are simply weaker than hoped as broad human-profile matches.

\begin{table}[H]
\centering
\small
\caption{Headline ARC-2 results for selected LLM and non-LLM profiles on the 110-pair primary subset.}
\begin{tabular}{p{2.45in}p{0.8in}p{0.9in}p{0.8in}}
\toprule
\textbf{Profile} & \textbf{Pair Acc.} & \textbf{Human $r$} & \textbf{H1 $p$} \\
\midrule
LLM average & 0.108 & 0.402 & 0.950 \\
Best-aligned LLM (Claude Opus 16k) & 0.164 & 0.439 & 0.506 \\
Best-score LLM (GPT-5-2 xhigh) & 0.491 & 0.276 & 0.052 \\
TRM 361957 pass@2 & 0.045 & 0.214 & 0.004 \\
TRM 651522 pass@2 & 0.109 & 0.113 & 0.0004 \\
VARC ARC-2 ViT pass@2 & 0.036 & 0.158 & 0.0004 \\
\bottomrule
\end{tabular}
\end{table}

\begin{table}[H]
\centering
\scriptsize
\caption{Threshold sensitivity on the ARC-2 public-eval overlap. Tightening the human-attempt filter changes the absolute values but not the broad ordering.}
\begin{tabular}{p{0.55in}p{0.55in}p{0.62in}p{0.62in}p{0.62in}p{0.62in}p{0.62in}p{0.62in}}
\toprule
\textbf{Min attempts} & \textbf{$n$ pairs} & \textbf{LLM avg} & \textbf{Opus 16k} & \textbf{GPT-5.2 xh} & \textbf{TRM 361957} & \textbf{TRM 651522} & \textbf{VARC ViT} \\
\midrule
2 & 161 & 0.345 & 0.348 & 0.256 & 0.164 & 0.090 & 0.091 \\
3 & 159 & 0.351 & 0.362 & 0.242 & 0.171 & 0.098 & 0.096 \\
5 & 129 & 0.367 & 0.403 & 0.231 & 0.199 & 0.095 & 0.142 \\
8 & 110 & 0.402 & 0.439 & 0.276 & 0.214 & 0.113 & 0.158 \\
\bottomrule
\end{tabular}
\end{table}

\begin{table}[H]
\centering
\small
\caption{Complementary rescue of human-easy and LLM-hard pairs. This is the compact reason the paper keeps a non-redundancy claim even while rejecting strong human-equivalence claims.}
\begin{tabular}{p{2.2in}p{0.8in}p{0.8in}p{0.9in}p{0.8in}}
\toprule
\textbf{System} & \textbf{Solved items} & \textbf{Rescued pairs} & \textbf{Random expectation} & \textbf{$p$} \\
\midrule
TRM 651522 pass@2 & 12 & 5 & 2.4 & 0.061 \\
VARC ARC-2 ViT pass@2 & 4 & 3 & 0.8 & 0.025 \\
TRM + VARC union & 13 & 6 & 2.6 & 0.022 \\
\bottomrule
\end{tabular}
\end{table}

\begin{figure}[H]
    \centering
    \includegraphics[width=0.92\textwidth]{fig02_threshold_sensitivity.png}
    \caption{Threshold-sensitivity figure from the non-LLM package. The ranking is stable enough that the paper can treat the disappointing non-LLM result as robust rather than a threshold artifact.}
\end{figure}

\begin{figure}[H]
    \centering
    \includegraphics[width=0.92\textwidth]{fig09_alignment_band.png}
    \caption{Selected ARC-2 system profiles against the human split-half benchmark. The visual point is simple: the LLM aggregate sits inside the human split-half band, while the current non-LLM systems do not.}
\end{figure}

\begin{figure}[H]
    \centering
    \includegraphics[width=0.92\textwidth]{fig10_fixed_accuracy_null.png}
    \caption{Observed human-alignment against the fixed-accuracy random-placement null. This is the key correction for score sparsity: low raw correlation can still be meaningful if it beats what random placement at the same number of wins would produce.}
\end{figure}

\begin{figure}[H]
    \centering
    \includegraphics[width=0.96\textwidth]{fig07_trm_trajectory.png}
    \caption{TRM training dynamics on the human-overlap subset. Score rises over training, but human-alignment peaks earlier and does not improve in lockstep.}
\end{figure}

\begin{figure}[H]
    \centering
    \includegraphics[width=0.92\textwidth]{fig11_residual_signal.png}
    \caption{Residual-signal view from the non-LLM package. The point is not that non-LLM systems broadly match humans, but that they retain some structure that is not exhausted by the LLM average.}
\end{figure}

\subsection{Why this section matters for the rest of the paper}

This section is where the repository's most disappointing result lives, and it is important precisely because the rest of the paper depends on not hiding it. Current non-LLM systems are not the clean human-like alternative to the LLM consensus on the available psychometric evidence. But they are also not vacuous. They contribute complementary successes, they separate benchmark score from human-like item placement, and they set up the architectural question explored later in the paper: why do these systems seem mechanistically promising even while their current human-alignment numbers remain modest?
